\scanpage{137}
% archive sheet number: 136

qui donne aussitôt les éléments
de structure

\[
\Delta_P,T_P,\varphi_P,\psi_P,\deleted{\unsure{\varphi}}\Delta_B,\psi_B,\varphi_B
\]

\uline{Th. récapitulatif} \textcircled{1} Les flèches
du diagramme de la page précédente
sont bijectives, le diagramme est
comm, et les \inserted{5} flèches qui ne sont
pas de source $\mathrm{Az}_P$ sont \deleted{\unsure{don}}compatibles
avec les structures de torseur
sous $(L\otimes L')^\vee$.

\textcircled{2} Pour qu'une structure \inserted{d'alg sur $P$} associée
à un twist en fasse une Alg.
d'Azumaya, il f et s. que
la forme $\Delta_B$ sur $B$ soit non
dégénérée, i.e que son discriminant
divisé

\[
\mathrm{disc}'(\Delta_B)\in(\det B)^{\otimes-2}=(L\otimes L')^{\otimes-2},
\]

soit inversible.

% slanted
\margin{Source \deleted{\unsure{et|cette}} et but ne dépendent que de $1_H,1_{H'},T_H,T_{H'}$ mais l'application $\delta$ \uline{fait} appel à $\Delta_H,\Delta_{H'}$}

\textcircled{3}

\diagram{\begin{tikzcd}[ampersand replacement=\&,column sep=huge]
\underline L_P\arrow[r,"{\delta_{P,B}}","{\substack{\text{discriminant}\\\text{divisé de}\\\text{la forme quad. sur}\\ B\ \text{associée}\\\text{à un twist}}}"']
\& (L\otimes L')^{\otimes-2}
\end{tikzcd}}

est une application \deleted{quadratique} pol. de degré $2$

